\section{Conclusion}
\label{sec:conclusion}

A stager's representation of a patient has two parts. One is a recording-wide
offset manufactured by the normalization statistics; the other is a
stage-conditional signature that describes how this patient's sleep stages
differ from the cohort prototypes. Marginal adaptation removes the first and,
as the within-stage probe shows, cannot target the second. Layerwise
adaptation may absorb part of it as a nonlinear side effect, but a
stage-conditional operation is required to remove the signature explicitly.
Stage-conditional shift normalization removes it with the model's own
posterior, in closed form, at the cost of one pass over the unlabeled night.

The experiments support three statements. The within-stage probe separates the
two components where the overall probe cannot, and it is the measurement a
subject-invariance claim should report. SCSN reduces that probe in every
setting, by between $0.19$ and $0.44$, and on the group-normalized checkpoints
it is the only stage-aware statistics-replacement method evaluated that does
so. Where it changes accuracy, it improves
it: significantly on the HMC backbones and on three of four group-normalized
settings, not measurably on Sleep-EDF EEGNet, and within noise of layerwise
AdaBN on both SleepCNN settings; the gains survive chance correction and class
balancing, and they come with a better worst-covered patient and a smaller
coverage dispersion.

The limitations define the next experiments. The method is evaluated on two
single-channel convolutional backbones, its estimate is only as good as the
posterior that indexes the stages, and it assumes the whole night is available.
A multi-channel or attention-based stager, a recalibrated posterior, and a
prefix-based online variant are the natural places to test whether the
decomposition --- rather than the particular estimator --- is what carries the
result.
